% ================================================================
% TOC ARTICLE PAGE-NUMBER EXAMPLE - Hebrew Academic Template
% ================================================================
% Regression example for v7.3.2:
%   Verifies that two-digit page numbers in the Table of Contents render
%   correctly (LTR) in the DEFAULT article mode.
%   Before v7.3.2 the TOC showed page 10 as "01", 12 as "21", 13 as "31", ...
%   because the book-mode \contentsline LTR fix was skipped in article mode.
%
% The document is intentionally padded with \newpage so the later sections
% land on two-digit pages, exercising the fix in the generated TOC.
%
% Compile: lualatex toc_article_pagenum_example.tex   (run twice for the TOC)
% Copyright (c) 2026 Dr. Segal Yoram. All rights reserved.
% ================================================================

\documentclass{hebrew-academic-template}   % article mode (default)

\usepackage{float}

\hebrewtitle{בדיקת מספרי עמודים בתוכן העניינים - מצב \en{article}}
\englishtitle{TOC Page-Number BiDi Test (article mode) \clsversion}
\hebrewauthor{ד"ר סגל יורם}
\hebrewversion{גרסה \en{\clsversion}}
\date{\textenglish{June 2026}}

\begin{document}

\maketitle
\tableofcontents
\newpage

% Each section is pushed onto a fresh page so that sections 8+ land on
% two-digit pages. The TOC must show those page numbers LTR (10, 11, 12...).

\hebrewsection{סעיף ראשון: \entoc{First Section}}
תוכן הסעיף הראשון. מספרים בתוך טקסט עברי כמו \num{100} ו-\percent{42} נשארים \en{LTR}.
\newpage

\hebrewsection{סעיף שני: \entoc{Second Section}}
תוכן הסעיף השני. מונחים באנגלית כמו \en{Machine Learning} משולבים בעברית.
\newpage

\hebrewsection{סעיף שלישי: \entoc{Third Section}}
תוכן הסעיף השלישי, כולל תת-סעיף.
\hebrewsubsection{תת-סעיף: \entoc{Subsection}}
מספר מורכב בתוכן העניינים, למשל \num{3.1}, צריך להופיע נכון.
\newpage

\hebrewsection{סעיף רביעי: \entoc{Fourth Section}}
תוכן הסעיף הרביעי.
\newpage

\hebrewsection{סעיף חמישי: \entoc{Fifth Section}}
תוכן הסעיף החמישי.
\newpage

\hebrewsection{סעיף שישי: \entoc{Sixth Section}}
תוכן הסעיף השישי.
\newpage

\hebrewsection{סעיף שביעי: \entoc{Seventh Section}}
תוכן הסעיף השביעי.
\newpage

\hebrewsection{סעיף שמיני: \entoc{Eighth Section}}
סעיף זה נמצא בעמוד דו-ספרתי. מספר העמוד שלו בתוכן העניינים חייב להיות \en{LTR}.
\newpage

\hebrewsection{סעיף תשיעי: \entoc{Ninth Section}}
תוכן הסעיף התשיעי.
\newpage

\hebrewsection{סעיף עשירי: \entoc{Tenth Section}}
תוכן הסעיף העשירי - עמוד דו-ספרתי.
\newpage

\hebrewsection{סעיף אחד-עשר: \entoc{Eleventh Section}}
תוכן הסעיף האחד-עשר.
\newpage

\hebrewsection{סעיף שנים-עשר: \entoc{Twelfth Section}}
תוכן הסעיף השנים-עשר. אם מספרי העמודים בתוכן העניינים מופיעים נכון
(\num{10}, \num{11}, \num{12} ולא \en{01}, \en{11}, \en{21}), התיקון של גרסה
\en{7.3.2} עובד כראוי.

\end{document}
